\begin{proof}[Proof of \cref{obj:lem:pm-reduced-slice-characterization}]
Write \(q=2(m-1)\) and keep the abbreviations \(x=1-u\), \(y=1+(m-1)u-mv\), \(z_{\mathrm{sp}}=1+(m-1)u+mv\) of the statement, so that
\[
y+z_{\mathrm{sp}}=2+2(m-1)u,
\qquad
z_{\mathrm{sp}}-y=2mv,
\qquad
qx+y+z_{\mathrm{sp}}=2m ,
\]
the last identity holding for all \(u,v\) because the \(u\)- and \(v\)-terms cancel.  Throughout, \((S_A,S_B)\) are the two community sums of \cref{obj:def:block-sum-handle}.


\emph{Part I: necessity.}  Suppose \(X(u,v)\in\mathcal C_m^{\pm}\), i.e.\ \(X(u,v)=X(P)=\mathbb E_P[ZZ^\top]\) for some \(P\in\mathcal P_m^{\mathrm{sym}}\).  Reading the entries of \(X(P)\) off the block pattern of \(X(u,v)\) gives \(\mathbb E[Z_iZ_j]=1\) if \(i=j\), \(=u\) if \(i\ne j\) lie in the same community, and \(=v\) if they lie in opposite communities.

\emph{(I.1) The three second-moment identities.}  Expanding \(S_A^2=\sum_{i,j\in A_m}Z_iZ_j\), there are \(m\) diagonal pairs and \(m(m-1)\) off-diagonal same-community pairs, and likewise for \(S_B^2\); expanding \(S_AS_B=\sum_{i\in A_m}\sum_{j\in B_m}Z_iZ_j\), all \(m^2\) pairs are cross-community.  Hence
\[
\mathbb E[S_A^2]=\mathbb E[S_B^2]=m+m(m-1)u,
\qquad
\mathbb E[S_AS_B]=m^2v .
\]
% lean: E_sumAr_sq % lean: E_sumBr_sq % lean: E_sumAr_sumBr

\emph{(I.2) Nonnegativity of \(y\) and \(z_{\mathrm{sp}}\).}  Combining the three identities,
\[
\mathbb E\bigl[(S_A-S_B)^2\bigr]
=2m+2m(m-1)u-2m^2v
=2m\,y,
\qquad
\mathbb E\bigl[(S_A+S_B)^2\bigr]
=2m+2m(m-1)u+2m^2v
=2m\,z_{\mathrm{sp}} .
\]
Both left-hand sides are expectations of squares, hence nonnegative, and \(m>0\); therefore \(y\ge0\) and \(z_{\mathrm{sp}}\ge0\). % lean: pm_slice_forward

\emph{(I.3) Nonnegativity of \(x\).}  Since \(|A_m|=m\ge2\), pick distinct \(i,j\in A_m\).  Because \(Z_i^2=Z_j^2=1\) pointwise,
\[
0\le\mathbb E\bigl[(Z_i-Z_j)^2\bigr]=2-2\,\mathbb E[Z_iZ_j]=2-2u,
\]
so \(u\le1\), i.e.\ \(x=1-u\ge0\). % lean: u_le_one

Together with the identity \(qx+y+z_{\mathrm{sp}}=2m\) recorded above, this gives the reduced-triangle conditions
\[
x\ge0,\qquad y\ge0,\qquad z_{\mathrm{sp}}\ge0,\qquad
2(m-1)x+y+z_{\mathrm{sp}}=2m .
\]

\emph{(I.4) The parity bound.}  If \(m\) is even then \(d_m=0\le y+z_{\mathrm{sp}}\) by (I.2), and there is nothing more to prove.  If \(m\) is odd, then \(S_A\) is a sum of \(m\) terms each equal to \(\pm1\) with \(m\) odd, so \(S_A\) is an odd integer; in particular \(S_A\ne0\) and, being an integer, \(S_A^2\ge1\).  The same holds for \(S_B\).  Taking expectations,
\[
\mathbb E[S_A^2]\ge1,
\qquad
\mathbb E[S_B^2]\ge1 .
\]
% lean: blockSumA_odd % lean: one_le_E_sumAr_sq % lean: one_le_E_sumBr_sq
On the other hand, by (I.1),
\[
\mathbb E[S_A^2]+\mathbb E[S_B^2]=2m+2m(m-1)u=m\,(y+z_{\mathrm{sp}}),
\]
so \(m(y+z_{\mathrm{sp}})\ge2\) and hence \(y+z_{\mathrm{sp}}\ge 2/m=d_m\). % lean: pm_slice_forward


\emph{Part II: sufficiency.}  Conversely, assume \(x,y,z_{\mathrm{sp}}\ge0\), \(qx+y+z_{\mathrm{sp}}=2m\) and \(y+z_{\mathrm{sp}}\ge d_m\).  We exhibit a law in \(\mathcal P_m^{\mathrm{sym}}\) whose second moment is \(X(u,v)\).

Two general facts are used.  First, the block parameterization is affine: for weights \(w_k\ge0\) with \(\sum_kw_k=1\),
\[
\sum_k w_k\,X(u_k,v_k)=X\Bigl(\sum_k w_ku_k,\ \sum_k w_kv_k\Bigr),
\]
because every diagonal entry of the left side is \(\sum_kw_k=1\) and the within- and across-community entries average separately. % lean: sum_blockSymMatrix
Second, if \(P=\sum_kw_kP_k\) is a convex mixture of laws \(P_k\in\mathcal P_m^{\mathrm{sym}}\), then \(P\in\mathcal P_m^{\mathrm{sym}}\) (sign symmetry and invariance under the two-block automorphism group are preserved by convex combinations) and, by linearity of expectation, \(X(P)=\sum_kw_kX(P_k)\). % lean: mixtureDesign_mem % lean: mixtureDesign_secondMoment

The building blocks are uniform laws on supports that are invariant under global sign reversal and under the two-block automorphism group; each such law lies in \(\mathcal P_m^{\mathrm{sym}}\), and its second moment is therefore block-symmetric, say \(X(u_\bullet,v_\bullet)\), with \(u_\bullet,v_\bullet\) determined through (I.1) by \(\mathbb E[S_A^2]\) and \(\mathbb E[S_AS_B]\):
\[
u_\bullet=\frac{\mathbb E[S_A^2]-m}{m(m-1)},
\qquad
v_\bullet=\frac{\mathbb E[S_AS_B]}{m^2}.
\]
% lean: secondMoment_blockSym_of_exchangeable

\emph{(II.1) Even \(m\): a three-vertex mixture.}  Take
\begin{itemize}
\item \(P_{\mathrm{cut}}\), uniform on the cut assignment (\(+1\) on \(A_m\), \(-1\) on \(B_m\)) and its global sign reversal, for which \(S_A^2=m^2\) and \(S_AS_B=-m^2\), hence \(X(P_{\mathrm{cut}})=X(1,-1)\); % lean: cutVDesign_secondMoment
\item \(P_{\mathrm{all}}\), uniform on the two constant assignments (all \(+1\), all \(-1\)), for which \(S_A^2=m^2\) and \(S_AS_B=m^2\), hence \(X(P_{\mathrm{all}})=X(1,1)\); % lean: allVDesign_secondMoment
\item \(P_{\mathrm{spr}}\), uniform on \(\{z:S_A(z)=S_B(z)=0\}\), a nonempty support when \(m\) is even, for which \(\mathbb E[S_A^2]=\mathbb E[S_AS_B]=0\) and hence \(X(P_{\mathrm{spr}})=X\!\left(-\tfrac{1}{m-1},0\right)\). % lean: spreadVDesign_secondMoment
\end{itemize}
Assign the weights
\[
w_{\mathrm{cut}}=\frac{y}{2m},
\qquad
w_{\mathrm{all}}=\frac{z_{\mathrm{sp}}}{2m},
\qquad
w_{\mathrm{spr}}=\frac{(m-1)x}{m}.
\]
They are nonnegative because \(x,y,z_{\mathrm{sp}}\ge0\) and \(m\ge2\), and they sum to one because
\[
\frac{y}{2m}+\frac{z_{\mathrm{sp}}}{2m}+\frac{(m-1)x}{m}
=\frac{y+z_{\mathrm{sp}}+2(m-1)x}{2m}
=\frac{2m}{2m}=1 .
\]
Their barycentre in the \((u,v)\) parameters is
\[
w_{\mathrm{cut}}\cdot1+w_{\mathrm{all}}\cdot1+w_{\mathrm{spr}}\cdot\Bigl(-\frac{1}{m-1}\Bigr)
=\frac{y+z_{\mathrm{sp}}}{2m}-\frac{x}{m}
=\frac{2m-2(m-1)x}{2m}-\frac{x}{m}
=1-x=u
\]
and
\[
w_{\mathrm{cut}}\cdot(-1)+w_{\mathrm{all}}\cdot1+w_{\mathrm{spr}}\cdot0
=\frac{z_{\mathrm{sp}}-y}{2m}=v .
\]
Hence the mixture \(P=w_{\mathrm{cut}}P_{\mathrm{cut}}+w_{\mathrm{all}}P_{\mathrm{all}}+w_{\mathrm{spr}}P_{\mathrm{spr}}\) lies in \(\mathcal P_m^{\mathrm{sym}}\) and has \(X(P)=X(u,v)\), so \(X(u,v)\in\mathcal C_m^{\pm}\). % lean: pm_slice_backward_even

\emph{(II.2) Odd \(m\): a four-vertex mixture.}  For odd \(m\) the support \(\{S_A=S_B=0\}\) is empty, so \(P_{\mathrm{spr}}\) is unavailable and the two parity vertices replace it:
\begin{itemize}
\item \(P^{\mathrm{par}}_{\mathrm{cut}}\), uniform on \(\{z:(S_A,S_B)=(1,-1)\text{ or }(-1,1)\}\) --- a nonempty support for odd \(m\) --- for which \(S_A^2=1\) and \(S_AS_B=-1\), hence
\[
X\bigl(P^{\mathrm{par}}_{\mathrm{cut}}\bigr)=X\Bigl(-\frac1m,\,-\frac1{m^2}\Bigr);
\]
% lean: pcutVDesign_secondMoment
\item \(P^{\mathrm{par}}_{\mathrm{all}}\), uniform on \(\{z:(S_A,S_B)=(1,1)\text{ or }(-1,-1)\}\), for which \(S_A^2=1\) and \(S_AS_B=1\), hence
\[
X\bigl(P^{\mathrm{par}}_{\mathrm{all}}\bigr)=X\Bigl(-\frac1m,\,\frac1{m^2}\Bigr).
\]
% lean: pallVDesign_secondMoment
\end{itemize}
In reduced coordinates these two points are \(\bigl(\tfrac{m+1}{m},\tfrac2m,0\bigr)\) and \(\bigl(\tfrac{m+1}{m},0,\tfrac2m\bigr)\), i.e.\ they sit on the parity boundary \(y+z_{\mathrm{sp}}=2/m\), whereas \(X(1,-1)\) and \(X(1,1)\) are the outer vertices \((0,2m,0)\) and \((0,0,2m)\).

Put \(s=y+z_{\mathrm{sp}}\) and
\[
\mu=\frac{y}{s},
\qquad
\lambda=\frac{s-2/m}{2m-2/m},
\]
which are well defined because \(s\ge2/m>0\) and \(2m-2/m>0\) for \(m\ge2\).  Then \(0\le\mu\le1\) since \(0\le y\le s\), and \(0\le\lambda\le1\) since \(2/m\le s\le 2m\), the upper bound coming from \(s=2m-qx\le2m\) with \(qx\ge0\).  Take the four weights
\[
\bigl(\lambda\mu,\ \lambda(1-\mu),\ (1-\lambda)\mu,\ (1-\lambda)(1-\mu)\bigr)
\]
on \(\bigl(P_{\mathrm{cut}},P_{\mathrm{all}},P^{\mathrm{par}}_{\mathrm{cut}},P^{\mathrm{par}}_{\mathrm{all}}\bigr)\); they are nonnegative and sum to one.  Their barycentre in the \(u\)-parameter is
\[
\lambda\cdot1+(1-\lambda)\Bigl(-\frac1m\Bigr)
=\frac{(m+1)\lambda-1}{m}
=u,
\]
where the last equality is the definition of \(\lambda\) rewritten through \(s=2m-2(m-1)(1-u)\), which gives \(\lambda=(1+mu)/(m+1)\).  For the \(v\)-parameter, the two \(\mu\)-pairs collapse to
\[
\lambda\mu(-1)+\lambda(1-\mu)(1)
+(1-\lambda)\mu\Bigl(-\frac1{m^2}\Bigr)+(1-\lambda)(1-\mu)\frac1{m^2}
=(1-2\mu)\Bigl(\lambda+\frac{1-\lambda}{m^2}\Bigr),
\]
and the two factors evaluate to
\[
1-2\mu=\frac{z_{\mathrm{sp}}-y}{s},
\qquad
\lambda+\frac{1-\lambda}{m^2}
=\frac{1}{m^2}+\lambda\,\frac{m^2-1}{m^2}
=\frac{s}{2m},
\]
the second using \(\lambda=(ms-2)/\bigl(2(m^2-1)\bigr)\).  Their product is \((z_{\mathrm{sp}}-y)/(2m)=v\).  Hence the four-point mixture lies in \(\mathcal P_m^{\mathrm{sym}}\) and has second moment \(X(u,v)\), so \(X(u,v)\in\mathcal C_m^{\pm}\). % lean: pm_slice_backward_odd


Parts I and II are the two implications of the asserted equivalence. % lean: pm_reduced_slice_characterization
\end{proof}
